%% Copyright (c) 2026, RTE (http://www.rte-france.com)
%% See AUTHORS.txt
%% All rights reserved.
%% This Source Code Form is subject to the terms of the Mozilla Public
%% License, v. 2.0. If a copy of the MPL was not distributed with this
%% file, you can obtain one at http://mozilla.org/MPL/2.0/.
%% SPDX-License-Identifier: MPL-2.0

\documentclass[a4paper, 12pt]{report}

%%  Copyright (c) 2015-2019, RTE (http://www.rte-france.com)
%%  See AUTHORS.txt
%%  All rights reserved.
%%  This Source Code Form is subject to the terms of the Mozilla Public
%%  License, v. 2.0. If a copy of the MPL was not distributed with this
%%  file, you can obtain one at http://mozilla.org/MPL/2.0/.
%%  SPDX-License-Identifier: MPL-2.0
%%
%%  This file is part of Dynawo, an hybrid C++/Modelica open source time domain
%%  simulation tool for power systems.


%%%%%%%%%%%%%%%%%%%%%%%%%%%%%%%%%%%%%%%%%%%
% Define text and document settings
%%%%%%%%%%%%%%%%%%%%%%%%%%%%%%%%%%%%%%%%%%%

\usepackage{lmodern} % Latin Modern fam­ily of fonts
\usepackage[english]{babel} % English

% Specify encoding
\usepackage[utf8]{inputenc} % Input
\usepackage[T1]{fontenc} % Output

% Document structure setup
\usepackage{titlesec} % To change chapter format
\setcounter{tocdepth}{3} % Add subsubsection in Content
\setcounter{secnumdepth}{3} % Add numbering for subsubsection
\setlength{\parindent}{0pt} % No paragraph indentation

% Avoid numbering starting at each chapter for figures
\usepackage{chngcntr}
\counterwithout{figure}{chapter}

% Change title format for chapter
\titleformat{\chapter}{\Huge\bf}{\thechapter}{20pt}{\Huge\bf}

% To add links on page number in Content and hide red rectangle on links
\usepackage[hidelinks, linktoc=all]{hyperref}
\usepackage[nottoc]{tocbibind} % To add biblio in table of content
\usepackage{textcomp} % For single quote
\usepackage{url} % Allow linebreaks in \url command
\usepackage{listings} % To add code samples

% Define typography
\usepackage{xspace}
\usepackage{dirtree}
\newcommand{\Dynawo}[0]{Dyna$\omega$o\xspace}

% Default listings parameters
\lstset
{
  aboveskip={1\baselineskip}, % A bit of space above
  backgroundcolor=\color{shadecolor}, % Choose the background color
  basicstyle={\ttfamily\footnotesize}, % Use font and smaller size \small \footnotesize
  breakatwhitespace=true, % Sets if automatic breaks should only happen at whitespace
  breaklines=true, % Sets automatic line breaking
  columns=fixed, % Nice spacing -> fixed / flexible
  mathescape=false, % Escape to latex false
  numbers=left, % Where to put the line-numbers
  numberstyle=\tiny\color{gray}, % The style that is used for the line-numbers
  showstringspaces=false, % Do not emphasize spaces in strings
  tabsize=4, % Number of spaces of a TAB
  texcl=false, % Activates or deactivates LaTeX comment lines
  upquote=true % Upright quotes
}

%%%%%%%%%%%%%%%%%%%%%%%%%%%%%%%%%%%%%%%%%%%
% Define plots settings
%%%%%%%%%%%%%%%%%%%%%%%%%%%%%%%%%%%%%%%%%%%

% Macro pack­age for cre­at­ing graph­ics
\usepackage{tikz}
\usepackage{subfigure}
\usepackage{float}

% Draws func­tion plots (based on pgf/tikz)
\usepackage{pgfplots}
\pgfplotsset{enlarge x limits=false, xlabel={\begin{small}$time$ (s)\end{small}}, height=0.6\textwidth, width=1\textwidth,
yticklabel style={text width={width("$-0.6$")},align=right}}
\pgfplotstableset{col sep=semicolon}

% Define colors
\usepackage{color}
\definecolor{blue}{rgb}{.3,.5,1}
\definecolor{deepblue}{rgb}{0,0,1}
\definecolor{darkblue}{rgb}{0,0,.4}
\definecolor{red}{rgb}{1,0,0}
\definecolor{darkred}{rgb}{.56,0,0}
\definecolor{pink}{rgb}{.933,0,.933}
\definecolor{purple}{rgb}{0.58,0,0.82}
\definecolor{green}{rgb}{0.133,0.545,0.133}
\definecolor{darkgreen}{rgb}{0,.4,0}
\definecolor{gray}{rgb}{.3,.3,.3}
\definecolor{darkgray}{rgb}{.2,.2,.2}
\definecolor{shadecolor}{gray}{0.925}

%%%%%%%%%%%%%%%%%%%%%%%%%%%%%%%%%%%%%%%%%%%
% Define blocks for simple network drawings
%%%%%%%%%%%%%%%%%%%%%%%%%%%%%%%%%%%%%%%%%%%

% Define blocks for newtorks drawings
\usepackage{amsmath} % Add math­e­mat­i­cal fea­tures
\usepackage{schemabloc} % Add block diagram library (french one)

%% Define infinite bus
\tikzset{infinite bus/.pic={
  code={
  \draw (0,0) circle (2) node[inner sep=0, outer sep=0] {{$\infty$}};
  \draw (2,0) --++ (2,0);
  }
  }
}

%% Define transformer
\tikzset{transfo/.pic={
  code={
  \draw (0,0) circle (2);
  \draw (2,0) circle (2);
  \draw (4,0) --++ (4,0);
  \draw (-2,0) --++ (-4,0);
  }
  }
}

%% Define generator
\tikzset{generator/.pic={
  code={
    \draw (0,0) circle (2);
    \draw (0,0) arc (0:180:0.5);
    \draw (0,0) arc (180:360:0.5);
    \draw (-2,0) --++ (-2,0);
  }
  }
}

%% Define generator controls
\tikzset{VR/.pic={
  code={
  \draw (0,0) circle (2) node[inner sep=0, outer sep=0] {{VR}};
  }
  }
}

%% Define SVarC
\tikzset{SVarC/.pic={
  code={
  \draw (0,0) circle (4) node[inner sep=0, outer sep=0] {{SVarC}};
  }
  }
}

\begin{document}
\title{DynaStream use examples}
\date\today
\maketitle
\tableofcontents

\chapter{Nordic test system}
\section{Test case description}
This example uses the Nordic network and its dynamic models to demonstrate an interactive simulation. The SIM solver runs from 0 to 3000 seconds with a coupling time step of one second. The scheduled line disconnection in \texttt{Events.par} is set to 10000000 seconds, outside this simulation; the example is intended for interactive changes instead.

\section{Models and inputs}
\texttt{Nordic.xiidm}, \texttt{Nordic.dyd} and their parameter files describe the network and dynamic models. \texttt{Events.dyd} adds updatable models for the status of line L4032-4042, load changes at 12\_1, and active and reactive power references at 11\_1. It also connects a set point to the active and reactive load changes at 61\_1, 62\_1 and 63\_1. \texttt{Nordic.crv} selects the curves to export.

\section{Running the simulation}
Use a Dynawo build with ZeroMQ support. From the Dynawo installation, run one of the following commands, replacing the example path as needed:
\begin{verbatim}
./dynawo.sh jobs /path/to/examples/DynaStream/Nordic/Nordic_internal_clock.jobs --interactive
./dynawo.sh jobs /path/to/examples/DynaStream/Nordic/Nordic_external_clock.jobs --interactive
\end{verbatim}
The internal clock advances at real-time speed by default. The external clock waits for a \texttt{step} control for each coupling step; without a client it does not progress. Both jobs bind the control channel on TCP port 5555 and the output channel on port 5556, so do not run them simultaneously without changing the endpoints.

\section{Interaction and results}
Install the Python dependencies from \texttt{rt-testing-tools/requirements.txt}. Start the subscriber in another terminal:
\begin{verbatim}
python rt-testing-tools/rt-testing-tools/main_zmq_subscriber.py
\end{verbatim}
Run this command from \texttt{examples/DynaStream}. It subscribes to all topics by default; \texttt{--topic} filters topics and \texttt{--dump\_dir} selects an existing directory for received dumps. The jobs stream curves, timeline and constraints as JSON; the external-clock job also enables dump streaming. Normal output files, including curves, timeline, constraints and final state, are written under \texttt{Nordic/outputs}.

The control helper can be used from a Python session in the tools directory:
\begin{verbatim}
from main_zmq_client import ControlSender
sender = ControlSender()
sender.send_step_trigger()  # repeat to advance the external clock
sender.send_dump_trigger()  # external-clock job streams dumps
sender.send_stop()
\end{verbatim}
The helper also exposes \texttt{send\_action} and \texttt{send\_load\_file}. Each request waits for an acknowledgement; start the simulation before sending requests. The supplied server is a separate control-channel test tool and cannot bind port 5555 while the simulation uses it.
\end{document}
